\documentclass[11pt,a4paper]{article}
\usepackage{fontspec,amsmath,amssymb,amsthm,microtype}
\setmainfont{lmroman10-regular.otf}[BoldFont=lmroman10-bold.otf,ItalicFont=lmroman10-italic.otf,BoldItalicFont=lmroman10-bolditalic.otf]
\usepackage{xcolor}
\usepackage[most]{tcolorbox}
\usepackage{mathrsfs}
\providecommand{\lt}{<}
\providecommand{\gt}{>}
\setmonofont{DejaVuSansMono.ttf}[BoldFont=DejaVuSansMono-Bold.ttf,ItalicFont=DejaVuSansMono-Oblique.ttf,BoldItalicFont=DejaVuSansMono-BoldOblique.ttf,Scale=0.86]
\newfontfamily\leanfallback{STIXTwoMath-Regular.otf}[Scale=0.86]
\definecolor{leanrule}{gray}{0.62}
\definecolor{leanback}{gray}{0.962}
\definecolor{leantitle}{gray}{0.22}
\definecolor{leanmuted}{gray}{0.38}
\newtcolorbox{leanbox}[1]{enhanced,breakable,colback=leanback,colframe=leanrule,
  boxrule=0.4pt,arc=0.6mm,left=2.4mm,right=2.4mm,top=1.8mm,bottom=1.8mm,
  before skip=0.8em,after skip=1em,fonttitle=\small\bfseries,coltitle=leantitle,
  colbacktitle=leanback,titlerule=0.3pt,title={#1}}
\newlength{\leancharwidth}
\newenvironment{leancode}{\par\smallskip\begingroup\ttfamily\footnotesize\raggedright
  \settowidth{\leancharwidth}{M}\setlength{\parskip}{0pt}\setlength{\parindent}{0pt}
  \frenchspacing}{\par\endgroup\smallskip}
\newcommand{\leanline}[2]{\hangindent=\dimexpr#1\leancharwidth+4\leancharwidth\relax
  \hangafter=1\noindent\hspace*{#1\leancharwidth}#2\par}
\newcommand{\leansource}[1]{\par{\footnotesize\color{leanmuted}#1}\par}
\newcommand{\leanlabel}[1]{\par\smallskip\noindent{\small\bfseries #1}\par}
\usepackage[margin=25mm]{geometry}
\usepackage[unicode,hidelinks]{hyperref}
\hypersetup{pdfauthor={LangevinLean Project}}
\setlength{\emergencystretch}{3em}
\setlength{\parskip}{0.4em}
\setlength{\parindent}{0pt}
\allowdisplaybreaks[2]
\raggedbottom
\title{\texorpdfstring{Strong-limit identities for actual observation compressions}{Strong-limit identities for actual observation compressions}}
\author{LangevinLean Project}
\date{}
\begin{document}\maketitle
Draft mathematical exposition; these notes may contain mistakes.


Common isometric source lifts do not change compression. Self-adjointness and composition identities pass through strong operator limits when the relevant approximants have uniform contraction bounds.


\section{\texorpdfstring{Isometric cancellation and strong-limit algebra}{Isometric cancellation and strong-limit algebra}}


\textbf{Claim.} Replacing both observation maps by their common isometric lift leaves their adjoint compression unchanged. A strong limit of self-adjoint operators is self-adjoint. If two strongly convergent operator families are contractions, their compositions converge strongly to the composition of their limits.



\textbf{Proof.} The lift cancels because its adjoint times itself is the identity. For self-adjointness, pass to the limit in the two inner-product pairings against fixed vectors. For composition, split the error into the first contraction applied to the second orbit error and the first-family error evaluated at the fixed limiting vector. Both terms tend to zero. Uniqueness of limits then identifies any separately known composition limit.

\begin{leanbox}{Lean: \texttt{compression\_\allowbreak{}isometricLift}, \texttt{strongLimit\_\allowbreak{}isSelfAdjoint}, \texttt{strongLimit\_\allowbreak{}comp} and 1 more}
\leansource{\texttt{compression\_\allowbreak{}isometricLift} \quad WassersteinMixing/\allowbreak{}SemigroupLimits.lean, lines 19–30}
\begin{leancode}
\leanline{0}{theorem\ compression\_isometricLift\ \{{\leanfallback 𝕜}\ A\ B\ C\ :\ Type*\}\ [RCLike\ {\leanfallback 𝕜}]}
\leanline{4}{[NormedAddCommGroup\ A]\ [InnerProductSpace\ {\leanfallback 𝕜}\ A]\ [CompleteSpace\ A]}
\leanline{4}{[NormedAddCommGroup\ B]\ [InnerProductSpace\ {\leanfallback 𝕜}\ B]\ [CompleteSpace\ B]}
\leanline{4}{[NormedAddCommGroup\ C]\ [InnerProductSpace\ {\leanfallback 𝕜}\ C]\ [CompleteSpace\ C]}
\leanline{4}{(initial\ current\ :\ A\ →ₗᵢ[{\leanfallback 𝕜}]\ B)\ (lift\ :\ B\ →ₗᵢ[{\leanfallback 𝕜}]\ C)\ :}
\leanline{4}{ObservableTransition.compression\ (lift.comp\ initial)\ (lift.comp\ current)\ =}
\leanline{6}{ObservableTransition.compression\ initial\ current}
\end{leancode}
\leansource{\texttt{strongLimit\_\allowbreak{}isSelfAdjoint} \quad WassersteinMixing/\allowbreak{}SemigroupLimits.lean, lines 36–49}
\begin{leancode}
\leanline{0}{theorem\ strongLimit\_isSelfAdjoint\ [l.NeBot]}
\leanline{4}{(A\ :\ I\ →\ H\ →L[{\leanfallback 𝕜}]\ H)\ (T\ :\ H\ →L[{\leanfallback 𝕜}]\ H)}
\leanline{4}{(hA\ :\ ∀\ i,\ IsSelfAdjoint\ (A\ i))}
\leanline{4}{(hlim\ :\ ∀\ f,\ Tendsto\ (fun\ i\ ↦\ A\ i\ f)\ l\ ({\leanfallback 𝓝}\ (T\ f)))\ :\ IsSelfAdjoint\ T}
\end{leancode}
\leansource{\texttt{strongLimit\_\allowbreak{}comp} \quad WassersteinMixing/\allowbreak{}SemigroupLimits.lean, lines 54–69}
\begin{leancode}
\leanline{0}{theorem\ strongLimit\_comp}
\leanline{4}{(A\ B\ :\ I\ →\ H\ →L[{\leanfallback 𝕜}]\ H)\ (T\ S\ :\ H\ →L[{\leanfallback 𝕜}]\ H)}
\leanline{4}{(hcontract\ :\ ∀\ i\ f,\ ‖A\ i\ f‖\ ≤\ ‖f‖)}
\leanline{4}{(hA\ :\ ∀\ f,\ Tendsto\ (fun\ i\ ↦\ A\ i\ f)\ l\ ({\leanfallback 𝓝}\ (T\ f)))}
\leanline{4}{(hB\ :\ ∀\ f,\ Tendsto\ (fun\ i\ ↦\ B\ i\ f)\ l\ ({\leanfallback 𝓝}\ (S\ f)))\ (f\ :\ H)\ :}
\leanline{4}{Tendsto\ (fun\ i\ ↦\ A\ i\ (B\ i\ f))\ l\ ({\leanfallback 𝓝}\ (T\ (S\ f)))}
\end{leancode}
\leansource{\texttt{strongLimit\_\allowbreak{}comp\_\allowbreak{}eq} \quad WassersteinMixing/\allowbreak{}SemigroupLimits.lean, lines 74–85}
\begin{leancode}
\leanline{0}{theorem\ strongLimit\_comp\_eq\ [l.NeBot]}
\leanline{4}{(A\ B\ C\ :\ I\ →\ H\ →L[{\leanfallback 𝕜}]\ H)\ (T\ S\ R\ :\ H\ →L[{\leanfallback 𝕜}]\ H)}
\leanline{4}{(hcontract\ :\ ∀\ i\ f,\ ‖A\ i\ f‖\ ≤\ ‖f‖)}
\leanline{4}{(hcomp\ :\ ∀\ i,\ C\ i\ =\ A\ i\ ∘L\ B\ i)}
\leanline{4}{(hA\ :\ ∀\ f,\ Tendsto\ (fun\ i\ ↦\ A\ i\ f)\ l\ ({\leanfallback 𝓝}\ (T\ f)))}
\leanline{4}{(hB\ :\ ∀\ f,\ Tendsto\ (fun\ i\ ↦\ B\ i\ f)\ l\ ({\leanfallback 𝓝}\ (S\ f)))}
\leanline{4}{(hC\ :\ ∀\ f,\ Tendsto\ (fun\ i\ ↦\ C\ i\ f)\ l\ ({\leanfallback 𝓝}\ (R\ f)))\ :\ R\ =\ T\ ∘L\ S}
\end{leancode}
\leanlabel{How the Lean proof follows the argument}
\texttt{compression\_\allowbreak{}isometric\allowbreak{}Lift} uses the exact adjoint composition identities. \texttt{strong\allowbreak{}Limit\_\allowbreak{}is\allowbreak{}Self\allowbreak{}Adjoint} passes symmetric pairings through limits. \texttt{strong\allowbreak{}Limit\_\allowbreak{}comp} uses the contraction norm bound and a two-term difference; \texttt{strong\allowbreak{}Limit\_\allowbreak{}comp\_\allowbreak{}eq} applies uniqueness of the resulting vector limits.
\end{leanbox}

\end{document}
